% GENERATED FILE. Edit canonical lessons, then run npm run generate:books.
% canonical-source-hash: fnv1a64:9eeb971bd7ce4b8a
% canonical-lessons: PT-C120-tornozelo, PT-C120-pulso, PT-C120-bochecha, PT-C120-labios, PT-C120-peito

\chapter{Ankle, Wrist, Cheek, Lips, Chest}
\label{ch:pt-ankle-wrist-cheek-lips-chest}
\hlchaptermodality{120}

\begin{chapteropening}
\textbf{By the end of this chapter:} \emph{I can say an ankle, a wrist, a cheek, lips and the chest.}

Complete the last lesson of chapter 120: I can say an ankle, a wrist, a cheek, lips and the chest.
\end{chapteropening}

\section[o tornozelo]{o tornozelo — an ankle}

\label{lesson:PT-C120-tornozelo}



\begin{quote}
Before the new one: say the Portuguese for a dolphin, then the Portuguese for an elbow.
\end{quote}

\subsection*{You'll want to know: o tornozelo}
\textbf{o tornozelo} — \textquotedblleft{}an ankle\textquotedblright{}.

\begin{cousinweb}
\textbf{tornozelo} is built on the word for a turn, with a small-making ending: the joint the foot turns on. EtymDB 2.1 links it to the verb \textquotedblleft{}to turn\textquotedblright{}, which $\leftarrow$ Latin \emph{tornare}, \textquotedblleft{}to turn on a lathe\textquotedblright{}.
\end{cousinweb}

\begin{grammarlens}[title={What you've built}]
One more word for this chapter.
\end{grammarlens}

\subsection*{Your turn}
\begin{itemize}
\raggedright
  \item \emph{Say it:} \emph{o tornozelo}
  \item \emph{Say it:} \emph{o tornozelo}, once more
  \item \emph{Say it:} say \emph{o tornozelo}
  \item \emph{Recall:} say the Portuguese for a dolphin, then the Portuguese for an elbow, then say \emph{o tornozelo} again
\end{itemize}

\begin{tcolorbox}[breakable,colback=teal!4,colframe=teal!35!black,title={Before you move on}]
What does o tornozelo mean? (\textquotedblleft{}An ankle\textquotedblright{}.) Say it once more.
\end{tcolorbox}
\section[o pulso]{o pulso — a wrist}

\label{lesson:PT-C120-pulso}



\begin{quote}
Before the new one: say the Portuguese for an elbow, then the Portuguese for an ankle.
\end{quote}

\subsection*{You'll want to know: o pulso}
\textbf{o pulso} — \textquotedblleft{}a wrist\textquotedblright{}.

\begin{cousinweb}
\textbf{pulso} $\leftarrow$ Latin \emph{pulsus}, \textquotedblleft{}a beat\textquotedblright{} (EtymDB 2.1), from \emph{pellere}, \textquotedblleft{}to drive\textquotedblright{} — a learned borrowing from Latin. The wrist is where the pulse beats.
\end{cousinweb}

\begin{grammarlens}[title={What you've built}]
One more word for this chapter.
\end{grammarlens}

\subsection*{Your turn}
\begin{itemize}
\raggedright
  \item \emph{Say it:} \emph{o pulso}
  \item \emph{Say it:} \emph{o pulso}, once more
  \item \emph{Say it:} say \emph{o pulso}
  \item \emph{Recall:} say the Portuguese for an elbow, then the Portuguese for an ankle, then say \emph{o pulso} again
\end{itemize}

\begin{tcolorbox}[breakable,colback=teal!4,colframe=teal!35!black,title={Before you move on}]
What does o pulso mean? (\textquotedblleft{}A wrist\textquotedblright{}.) Say it once more.
\end{tcolorbox}
\section[a bochecha]{a bochecha — a cheek}

\label{lesson:PT-C120-bochecha}



\begin{quote}
Before the new one: say the Portuguese for an ankle, then the Portuguese for a wrist.
\end{quote}

\subsection*{You'll want to know: a bochecha}
\textbf{a bochecha} — \textquotedblleft{}a cheek\textquotedblright{}.

\begin{cousinweb}
\textbf{bochecha} is probably built on \emph{boca}, \textquotedblleft{}mouth\textquotedblright{} ($\leftarrow$ Latin \emph{bucca}, \textquotedblleft{}cheek\textquotedblright{}); its exact formation is uncertain.
\end{cousinweb}

\begin{grammarlens}[title={What you've built}]
One more word for this chapter.
\end{grammarlens}

\subsection*{Your turn}
\begin{itemize}
\raggedright
  \item \emph{Say it:} \emph{a bochecha}
  \item \emph{Say it:} \emph{a bochecha}, once more
  \item \emph{Say it:} say \emph{a bochecha}
  \item \emph{Recall:} say the Portuguese for an ankle, then the Portuguese for a wrist, then say \emph{a bochecha} again
\end{itemize}

\begin{tcolorbox}[breakable,colback=teal!4,colframe=teal!35!black,title={Before you move on}]
What does a bochecha mean? (\textquotedblleft{}A cheek\textquotedblright{}.) Say it once more.
\end{tcolorbox}
\section[os lábios]{os lábios — lips}

\label{lesson:PT-C120-labios}



\begin{quote}
Before the new one: say the Portuguese for a wrist, then the Portuguese for a cheek.
\end{quote}

\subsection*{You'll want to know: os lábios}
\textbf{os lábios} — \textquotedblleft{}lips\textquotedblright{}.

\begin{cousinweb}
\textbf{lábios} $\leftarrow$ Latin \emph{labium}, \textquotedblleft{}lip\textquotedblright{} (Lewis \& Short) — a learned borrowing from Latin.
\end{cousinweb}

\begin{grammarlens}[title={What you've built}]
One more word for this chapter.
\end{grammarlens}

\subsection*{Your turn}
\begin{itemize}
\raggedright
  \item \emph{Say it:} \emph{os lábios}
  \item \emph{Say it:} \emph{os lábios}, once more
  \item \emph{Say it:} say \emph{os lábios}
  \item \emph{Recall:} say the Portuguese for a wrist, then the Portuguese for a cheek, then say \emph{os lábios} again
\end{itemize}

\begin{tcolorbox}[breakable,colback=teal!4,colframe=teal!35!black,title={Before you move on}]
What does os lábios mean? (\textquotedblleft{}Lips\textquotedblright{}.) Say it once more.
\end{tcolorbox}
\section[o peito]{o peito — the chest}

\label{lesson:PT-C120-peito}



\begin{quote}
Before the new one: say the Portuguese for a cheek, then the Portuguese for lips.
\end{quote}

\subsection*{You'll want to know: o peito}
\textbf{o peito} — \textquotedblleft{}the chest\textquotedblright{}.

\begin{cousinweb}
\textbf{peito} $\leftarrow$ Latin \emph{pectus}, \textquotedblleft{}breast\textquotedblright{} (EtymDB 2.1) — inherited. Latin \emph{-ct-} became \emph{-it-}.
\end{cousinweb}

\begin{grammarlens}[title={What you've built}]
One more word for this chapter.
\end{grammarlens}

\subsection*{Your turn}
\begin{itemize}
\raggedright
  \item \emph{Say it:} \emph{o peito}
  \item \emph{Say it:} \emph{o peito}, once more
  \item \emph{Say it:} say \emph{o peito}
  \item \emph{Recall:} say the Portuguese for a cheek, then the Portuguese for lips, then say \emph{o peito} again
\end{itemize}

\begin{tcolorbox}[breakable,colback=teal!4,colframe=teal!35!black,title={Before you move on}]
What does o peito mean? (\textquotedblleft{}The chest\textquotedblright{}.) Say it once more.
\end{tcolorbox}
